\documentclass{article}
\usepackage[final,nonatbib]{neurips_2024}
\usepackage[T1]{fontenc}
\usepackage{amsmath,amssymb,booktabs,graphicx,url}
\usepackage[hidelinks]{hyperref}
\makeatletter
\renewcommand{\@noticestring}{ECE 285: Homework 2.}
\makeatother
\newcommand{\TrainCount}{75,694}
\newcommand{\RawCount}{81,444}
\newcommand{\BadCount}{0}
\newcommand{\DuplicateCount}{750}
\newcommand{\UniqueCount}{80,694}
\newcommand{\StepCount}{5,910}
\newcommand{\FinalFID}{95.30}
\newcommand{\FinalIS}{2.892}
\newcommand{\RealIS}{7.378}
\newcommand{\InitialFID}{471.87}
\newcommand{\InitialIS}{1.001}
\newcommand{\FirstGLoss}{6.15}
\newcommand{\FirstDLoss}{0.58}
\newcommand{\FinalGLoss}{3.30}
\newcommand{\FinalDLoss}{0.61}
\newcommand{\MetricRows}{
Untrained & 471.87 & 1.001 & $1.001 \pm 0.00003$ \\
Epoch 1 & 333.56 & 1.676 & $1.672 \pm 0.032$ \\
Epoch 5 & 151.99 & 2.280 & $2.266 \pm 0.038$ \\
Epoch 10 & 95.30 & 2.892 & $2.861 \pm 0.093$ \\
}
\title{Generating WikiArt Paintings with a DCGAN}
\author{ECE 285 HW2 \quad Code: \url{https://github.com/alextung2/ECE-285-HW2}}
\begin{document}
\maketitle
\begin{abstract}
A DCGAN trained from scratch on \TrainCount\ WikiArt images generates
$64\times64$ RGB paintings. After ten epochs, evaluation against 5,000
held-out images gives FID \FinalFID\ and Inception Score \FinalIS\ under
the stated torchvision-Inception protocol. Samples exhibit coarse structure,
blur, and residual grid artifacts.
\end{abstract}

\section{Data and model}
The assigned Kaggle archive~[1] contains \RawCount\ images, all successfully
decoded. Each is converted to RGB, bilinearly resized to a 64-pixel shorter
side, and center-cropped to $64\times64$. Removing \DuplicateCount\ identical
crops (SHA-256) leaves \UniqueCount\ images. A seeded permutation reserves
5,000 for evaluation and \TrainCount\ for training. Pixels are scaled to
$[-1,1]$; style labels are ignored. The split is image-level, with no
artist-level or perceptual-near-duplicate separation.

The DCGAN~[2] generator maps $z\sim\mathcal N(0,I_{100})$ through five
$4\times4$ transposed convolutions: channels $512,256,128,64,3$, spatial sizes
$4,8,16,32,64$. Hidden layers use batch normalization and ReLU; the output
uses tanh. The discriminator uses $4\times4$ convolutions with channels
$64,128,256,512,1$, LeakyReLU($0.2$), batch normalization in the three
intermediate blocks, and a final logit. Generator/discriminator parameter
counts are 3,576,704/2,765,568.

\section{Optimization and training dynamics}
With $D$ denoting the sigmoid of the discriminator logit, the objectives are
\begin{align*}
L_D&=-\mathbb E_x\log D(x)-\mathbb E_z\log[1-D(G(z))],&
L_G&=-\mathbb E_z\log D(G(z)).
\end{align*}
\noindent Stable BCE-with-logits implements both losses. Each step updates $D$
using real and detached fake images, then $G$ using the updated $D$.
Both use Adam, learning rate $2\times10^{-4}$, $\beta=(0.5,0.999)$, and batch
size 128 (incomplete batches dropped). Convolution weights start at
$\mathcal N(0,0.02^2)$; batch-norm scales at $\mathcal N(1,0.02^2)$ and biases
at zero. Seed 285, float32, and an RTX 3080 Ti Laptop GPU are used for
\StepCount\ updates. No pretrained generator or augmentation is used.

\begin{figure}[!ht]
\centering
\includegraphics[width=\linewidth]{loss_curve.pdf}
\caption{Actual per-step BCE losses (faint) and trailing 100-step means (solid).
The discriminator loss sums real and fake terms. Adversarial losses need not
decrease monotonically and do not establish perceptual quality.}
\end{figure}

\newpage
\section{Generated images and quantitative evaluation}
\begin{figure}[h]
\centering
\includegraphics[width=\linewidth]{comparison.pdf}
\caption{First 16 held-out images and first 16 fixed-noise outputs.
Generated panels share latent vectors; examples are not selected by visual quality.}
\end{figure}

Evaluation uses Ignite 0.5.3 FID and InceptionScore, adapting~[3]. Each checkpoint
uses 5,000 Gaussian latents (seed 1285, reset per checkpoint) and the same 5,000
held-out images. Both inputs are mapped to $[0,1]$, identically rounded to 8-bit
RGB, PIL-bilinearly resized to $299\times299$, and ImageNet-normalized.
Frozen torchvision Inception-v3 (\texttt{IMAGENET1K\_V1},
\texttt{transform\_input=True}) supplies 2,048-dimensional average-pool features
for FID, explicitly replacing Ignite's default 1,000 probabilities~[4].
IS uses softmax probabilities. Lower FID and higher IS are preferable.

\begin{table}[h]
\centering
\caption{Measured results under one fixed protocol. Whole-set IS is Ignite's
direct output; split IS is mean $\pm$ standard deviation over ten groups of 500.}
\small
\begin{tabular}{lrrr}
\toprule
Checkpoint & FID $\downarrow$ & IS, whole set $\uparrow$ & IS, ten splits $\uparrow$\\
\midrule
\MetricRows
\bottomrule
\end{tabular}
\end{table}

\paragraph{Observed behavior.}
Final outputs show coarse landscape-like arrangements, with blur and grid
artifacts. Mean generator loss falls from \FirstGLoss\ (epoch 1) to
\FinalGLoss\ (epoch 10); discriminator loss changes from \FirstDLoss\ to
\FinalDLoss. FID improves from \InitialFID\ to \FinalFID\ and IS from
\InitialIS\ to \FinalIS. These trends support learning, but do not establish
convergence or complete style coverage.


\paragraph{Limitations.}
Torchvision-based FID is not directly comparable to TensorFlow/clean-fid
benchmarks. Finite-sample bias and a single training seed limit inference.
Split-IS deviation measures sample-group variation, not training-run uncertainty.
ImageNet classes weakly reflect artistic quality (real-set IS: \RealIS).
Cropping removes composition and $64\times64$ limits detail. Losses do not
establish absence of memorization or complete mode coverage. More training,
multiple seeds, and perceptual duplicate checks would strengthen the study.
Code, splits, logs, and checkpoints are saved; cuDNN benchmarking prevents
guaranteed bit-identical CUDA reproduction.

\vspace{-3pt}
\section*{References}
\vspace{-3pt}
{\footnotesize\raggedright
\noindent [1] steubk. \emph{WikiArt}. Kaggle.
\url{https://www.kaggle.com/datasets/steubk/wikiart}.\\
\noindent [2] A. Radford, L. Metz, and S. Chintala. \emph{Unsupervised Representation
Learning with Deep Convolutional Generative Adversarial Networks}. ICLR, 2016.
\url{https://arxiv.org/abs/1511.06434}.\\
\noindent [3] PyTorch-Ignite. \emph{GAN evaluation using FID and IS}.
\url{https://pytorch-ignite.ai/blog/gan-evaluation-with-fid-and-is/}.\\
\noindent [4] PyTorch-Ignite. \emph{FID metric documentation}.
\url{https://docs.pytorch.org/ignite/generated/ignite.metrics.FID.html}.\par
}
\end{document}
